\subsectionaddtolist{رمزنگاری RSA}

\subsubsection*{آ) محاسبه کلید عمومی $(n, e)$ و کلید خصوصی $(n, d)$:}

ابتدا مقدار $n$ را محاسبه می‌کنیم:
$$n = p \times q = 7 \times 13 = 91$$

سپس مقدار $\phi(n)$ را به دست می‌آوریم:
$$\phi(n) = (p-1)(q-1) = (7-1)(13-1) = 6 \times 12 = 72$$

برای یافتن کلید خصوصی $d$، باید معکوس ضربی $e=7$ را به پیمانه $\phi(n)=72$ با استفاده از {الگوریتم اقلیدسی تعمیم‌یافته} پیدا کنیم. یعنی حل معادله هم‌نهشتی:
$$ 7d \equiv 1 \pmod{72}$$

مراحل الگوریتم تقسیم اقلیدسی (رو به جلو):
$$72 = 10 \times 7 + 2$$
$$7 = 3 \times 2 + 1$$

حالا مراحل را به صورت بازگشتی (رو به عقب) برای یافتن ترکیب خطی بنویسیم:
$$1 = 7 - 3 \times 2$$
جایگذاری $2 = 72 - 10 \times 7$:
$$1 = 7 - 3 \times (72 - 10 \times 7)$$
$$1 = 7 - 3 \times 72 + 30 \times 7$$
$$1 = 31 \times 7 - 3 \times 72$$

بنابراین:
$$31 \times 7 \equiv 1 \pmod{72} \implies d = 31$$

{کلید عمومی:} $(n, e) = (91, 7)$

{کلید خصوصی:} $(n, d) = (91, 31)$


\subsubsection*{ب) رمزگذاری پیام $M = 82$ و محاسبه متن رمزنشده $C$:}

فرمول رمزگذاری:
$$C \equiv M^e \pmod{n} \implies C \equiv 82^7 \pmod{91}$$

ابتدا پایه را ساده می‌کنیم: $82 \equiv -9 \pmod{91}$. پس:
$$C \equiv (-9)^7 \pmod{91}$$

$$(-9)^2 = 81 \equiv -10 \pmod{91}$$
$$(-9)^4 \equiv (-10)^2 = 100 \equiv 9 \pmod{91}$$
$$(-9)^6 = (-9)^4 \times (-9)^2 \equiv 9 \times (-10) = -90 \equiv 1 \pmod{91}$$

در نهایت:
$$C \equiv (-9)^7 = (-9)^6 \times (-9)^1 \equiv 1 \times (-9) = -9 \equiv 82 \pmod{91}$$

پس متن رمزگذاری شده برابر است با: \textbf{$C = 82$}

\subsubsection*{ج) رمزگشایی متن $C$ و بازگشت به پیام اصلی $M$:}

فرمول رمزگشایی:
$$M = C^d \pmod{n} \implies M = 82^{31} \pmod{91}$$

با توجه به این که $82 \equiv -9 \pmod{91}$، داریم:
$$M \equiv (-9)^{31} \pmod{91}$$

با استفاده از محاسبات بخش قبل می‌دانیم که $(-9)^6 \equiv 1 \pmod{91}$. بنابراین توان را به پیمانه ۶ ساده می‌کنیم:
$$31 = 5 \times 6 + 1$$
$$(-9)^{31} = ((-9)^6)^5 \times (-9)^1 \equiv (1)^5 \times (-9) = -9 \equiv 82 \pmod{91}$$

پس پیام رمزگشایی شده برابر با \textbf{$M = 82$} است که با پیام اولیه کاملا مطابقت دارد.

\subsectionaddtolist{تبادل کلید Diffie-Hellman}

\subsubsection*{آ) یافتن کلید خصوصی آلیس ($X_A$) در صورت داشتن کلید عمومی $Y_A = 9$:}

فرمول کلید عمومی در پروتکل دیفی-هلمن عبارت است از:
$$Y_A = \alpha^{X_A} \pmod{q} \implies 9 = 2^{X_A} \pmod{11}$$

با آزمون مقادیر مختلف برای $X_A$ (از $1$ تا $10$):
\begin{itemize}
    \item $2^1 = 2 \pmod{11}$
    \item $2^2 = 4 \pmod{11}$
    \item $2^3 = 8 \pmod{11}$
    \item $2^4 = 16 \equiv 5 \pmod{11}$
    \item $2^5 = 32 \equiv 10 \pmod{11}$
    \item $2^6 = 64 \equiv 9 \pmod{11}$
\end{itemize}

بنابراین، کلید خصوصی آلیس برابر است با: \textbf{$X_A = 6$}

مسئله پیدا کردن $X_A$ از روی $Y_A$، $\alpha$ و $q$ به نام {مسئله لگاریتم گسسته} شناخته می‌شود. برای اعداد اول بسیار بزرگ ($q$ با چندصد رقم)، هیچ الگوریتم کلاسیک با زمان چندجمله‌ای برای حل این مسئله در زمان معقول وجود ندارد و امنیت این پروتکل نیز دقیقا بر پایه همین دشواری استوار است.

\subsubsection*{ب) محاسبه کلید مشترک $K_{AB}$ در صورت داشتن کلید عمومی باب $Y_B = 3$:}

هر دو طرف می‌توانند به طور مستقل کلید مشترک را محاسبه کنند:

\begin{enumerate}
    \item {محاسبه توسط آلیس (با استفاده از کلید عمومی باب $Y_B$ و کلید خصوصی خود $X_A = 6$):}
    $$K_{AB} = Y_B^{X_A} \pmod{q} = 3^6 \pmod{11}$$
    $$3^2 = 9 \equiv -2 \pmod{11}$$
    $$3^4 \equiv (-2)^2 = 4 \pmod{11}$$
    $$3^6 = 3^4 \times 3^2 \equiv 4 \times 9 = 36 \equiv 3 \pmod{11}$$
    بنابراین کلید مشترک از نظر آلیس $3$ است.

    \item {محاسبه توسط باب (ابتدا نیاز به یافتن کلید خصوصی خود $X_B$ دارد):}
    $$Y_B = \alpha^{X_B} \pmod{q} \implies 3 = 2^{X_B} \pmod{11}$$
    با ادامه بررسی توان‌های ۲ از بخش قبل:
    \begin{itemize}
        \item $2^7 = 2^6 \times 2 \equiv 9 \times 2 = 18 \equiv 7 \pmod{11}$
        \item $2^8 = 2^7 \times 2 \equiv 7 \times 2 = 14 \equiv 3 \pmod{11}$
    \end{itemize}
    پس کلید خصوصی باب $X_B = 8$ است. حالا باب کلید مشترک را با کلید عمومی آلیس ($Y_A = 9$) محاسبه می‌کند:
    $$K_{AB} = Y_A^{X_B} \pmod{q} = 9^8 \pmod{11}$$
    از آنجا که $9 \equiv -2 \pmod{11}$:
    $$9^8 \equiv (-2)^8 = 256 \pmod{11}$$
    با تقسیم ۲۵۶ بر ۱۱ داریم: $256 = 23 \times 11 + 3$. پس $256 \equiv 3 \pmod{11}$.
\end{enumerate}

مشاهده می‌شود که هر دو طرف به طور مستقل به کلید مشترک یکسان \textbf{$K_{AB} = 3$} رسیدند.